\subsection{互译}

设 K 是（交换）域，f 是 K 的位，v 是 K 上的赋值，A 是 K 的赋值环。若 A 既是 f 的环又是 v 的环，则称 A、f 与 v 相伴。由第 1 节与 § 2 第 3 节，三个对象 A、f 与 v 中的每一个都确定另外两个（对位与赋值而言，确定到相差一个等价）。特别地，有下列等价关系：

\[
\begin{array}{l l l l} x \in \mathrm{A} & \Leftrightarrow f (x) \neq \infty & \Leftrightarrow v (x) \geqslant 0 \\ x \in \mathfrak {m} (\mathrm{A}) & \Leftrightarrow f (x) = 0 & \Leftrightarrow v (x) > 0 \\ x \in \mathrm{A} - \mathfrak {m} (\mathrm{A}) = \mathrm{U} (\mathrm{A}) \Leftrightarrow f (x) \neq 0 \text {and} f (x) \neq \infty \Leftrightarrow v (x) = 0 \\ x \in \mathrm{K} - \mathrm{A} & \Leftrightarrow f (x) = \infty & \Leftrightarrow v (x) <   0 \end{array}
\]

关于赋值环、位或赋值的每个结果都可以翻译成关于另外两个概念的结果。于是 § 2 第 4 节命题 4 给出：

命题 5. 设 K 是域，v 是 K 上的赋值，K' 是 K 的扩张。存在 K' 上的赋值 v'，它在 K 上的限制等价于 v。

设 \(\Gamma_v\) 与 \(\Gamma_{v'}\) 是 \(v\) 与 \(v'\) 的阶群。由于 \(v'\) 在 K 上的限制等价于 \(v\)，存在同构 \(\lambda\)，它把 \(\Gamma_v\) 映到 \(\Gamma_{v'}\) 的一个子群上，使在 K 上 \(v' = \lambda \circ v\)。若把 \(\Gamma_v\) 等同于 \(\lambda(\Gamma_v)\)（借助 \(\lambda\)），则可见 \(v'\) 延拓 \(v\)。

注意 \(\Gamma_{v'}\) 一般异于 \(\lambda(\Gamma_v)\)，且 \(v'\) 的等价类不必唯一。我们将在 § 8 中回到这一点。

翻译 § 1 第 3 节定理 3（或 § 2 第 5 节命题 6），我们得到：

命题 6. 设 K 是域，A 是 K 的子环，x 是 K 的元素。为使 x 在 A 上整，必须且只需 K 的每个在 A 上为正的赋值都在 x 处为正。

从现在起，我们将把进行与上述类似的翻译的工作一般留给读者。

